% from main.tex
%%%%%%%%%%  scan idx 121  |  folio 108  %%%%%%%%%%
(i) \emph{\(\mathcal C_1\) est localement quasi-relevable dans \(\mathcal C\) (1.2) et localement quasi-stable par noyau d'épimorphisme (1.2).}

(ii) \emph{Soient \(S\in\Ob S\) et \(F\in\Ob\Dpx(\mathcal C_S)\). Pour que \(F\) soit \(n\)-pseudo-cohérent (resp. pseudo-cohérent) relativement à \(\mathcal C_1\), il faut et il suffit qu'il le soit relativement à \(\mathcal C_0\).}

\emph{Preuve.} (i) Soit \(E\to F\) un épimorphisme de \(\mathcal C_S\), avec \(F\in\Ob\mathcal C_{1S}\). Comme \(F\) est \(\mathcal C_0\)-pseudo-cohérent (donc a fortiori de \(\mathcal C_0\)-type fini), il existe (quitte à localiser) un épimorphisme \(E'\to F\) avec \(E'\in\Ob\mathcal C_{0S}\). Utilisant la quasi-relevabilité de \(\mathcal C_0\) dans \(\mathcal C\), on peut supposer (quitte à localiser à nouveau) que celui-ci se factorise à travers \(E\), d'où la quasi-relevabilité de \(\mathcal C_1\). D'autre part, si \(E\in\Ob\mathcal C_S\) admet une résolution droite finie par des objets de \(\mathcal C_1\), \(E\) est \(\mathcal C_0\)-pseudo-cohérent en vertu de (2.11.a)), donc a fortiori de \(\mathcal C_1\)-type fini.

(ii) Il suffit de prouver la nécessité. On peut pour cela appliquer, au choix, (2.11.a)) ou (2.13) avec \(S'=S\), \(\mathcal C'=\mathcal C\), \(f=\text{Identité}\), \(u=\text{Identité}\), \(\mathcal C'_0\) et \(\mathcal C_0\) remplacées respectivement par \(\mathcal C_0\) et \(\mathcal C_1\)).

\textbf{Scholie 2.14.1.} Il existe donc une plus grande sous-\(S\)-catégorie additive strictement pleine de \(\mathcal C\) donnant la même notion de \(n\)-pseudo-cohérence (resp. pseudo-cohérence) que \(\mathcal C_0\), à savoir la catégorie formée des objets \(\mathcal C_0\)-pseudo-cohérents. Cette dernière est localement stable par noyau d'épimorphisme.

\textbf{Exemple 2.15.} Soient \((S,A)\) un topos annelé et \(\mathcal C,\mathcal C_0\) les \(S\)-catégories envisagées en (1.3.d)). La catégorie fibrée \(\Dpx(\mathcal C)\) sera souvent notée \(\Dpx(S,A)\), ou \(\Dpx({}_A S)\) (pour
